% ============================================================
%  teacher_guide.tex — Inside a Smartphone: The Tiny World in Your Hand
%  Parent / Teacher Guide: discussion questions, extension
%  activities, and curriculum alignment per chapter.
% ============================================================

\namedchapter{Parent \& Teacher Guide}
\addcontentsline{toc}{chapter}{Parent \& Teacher Guide}

\begin{tcolorbox}[enhanced,colback=SunYellow!12,colframe=TechPurple!55,arc=6pt,boxrule=1pt,left=10pt,right=10pt,top=8pt,bottom=8pt]
\textbf{Vocabulary scaffolding note (Chapters 2--4):} The terms transistors, gigahertz, triangulation, capacitive, operating system, and biometrics may need additional support. Consider pronunciation guides or visual vocabulary icons.\\
\textbf{Chapter 10 engagement note:} The tone shifts toward instruction in this chapter. Add scenario-based examples or mini decision-making discussions to keep narrative engagement strong.
\end{tcolorbox}

\begin{tcolorbox}[enhanced,colback=LightGray,colframe=TechPurple!50,arc=6pt,boxrule=1pt,
  left=10pt,right=10pt,top=10pt,bottom=10pt]
\textbf{How to use this guide:}
Each chapter entry includes discussion questions to explore together,
extension activities for deeper learning, and notes on curriculum areas supported.
Read one chapter at a time and allow space for conversation between chapters.
This book builds on concepts from \emph{The Secret Life of the Internet} (Book 1)
but can also be read as a standalone.
\end{tcolorbox}

\vspace{12pt}

% ============================================================
\section*{Chapter 1: What Is a Smartphone?}
% ============================================================

\subsection*{Discussion Questions}
\begin{enumerate}
  \item Name five different things you can do with a smartphone. Would you have needed five different devices 30 years ago?
  \item How do you think life is different for people who do not have a smartphone?
  \item A smartphone is described as a tiny computer. What makes it similar to a laptop? What makes it different?
\end{enumerate}

\subsection*{Extension Activities}
\begin{itemize}
  \item \textbf{Feature audit:} List every feature and sensor your household phone has and categorise them (communication, navigation, camera, etc.).
  \item \textbf{Timeline:} Research the history of mobile phones from the 1970s brick phone to today's smartphones. Create a visual timeline.
  \item \textbf{Comparison chart:} Compare a smartphone to a 1990s desktop computer — processing speed, storage, camera, portability.
\end{itemize}

\subsection*{Curriculum Links}
Science \& Technology, History (technological change), PSHE (digital citizenship).

\sectionrule

% ============================================================
\section*{Chapter 2: Meet the Brain: The Processor}
% ============================================================

\subsection*{Discussion Questions}
\begin{enumerate}
  \item Why do you think processor speed is measured in \emph{billions} of instructions per second? Can you imagine what one billion of anything looks like?
  \item Why do some apps make your phone warm? Where does that heat come from?
  \item What would happen if a phone had a very slow processor?
\end{enumerate}

\subsection*{Extension Activities}
\begin{itemize}
  \item \textbf{Instruction game:} Give a child a very detailed set of instructions for a simple task (e.g., making a sandwich). Time them. Discuss how computers must follow instructions even more precisely.
  \item \textbf{Research:} Find out how many transistors are in a modern smartphone processor versus a 1970s CPU.
  \item \textbf{Human CPU:} Act out a processor — one person gives instructions, another executes them, a third records the result.
\end{itemize}

\subsection*{Curriculum Links}
Mathematics (large numbers, scale), Computing (algorithms, logical thinking), Science (heat energy).

\sectionrule

% ============================================================
\section*{Chapter 3: Memory and Storage}
% ============================================================

\subsection*{Discussion Questions}
\begin{enumerate}
  \item What is the difference between RAM and storage in your own words?
  \item Why do you think photos and videos take up more storage than text messages?
  \item What would happen if you deleted all your apps — would the phone still work?
\end{enumerate}

\subsection*{Extension Activities}
\begin{itemize}
  \item \textbf{Try It:} Check the storage on a device. How much is used? What is using the most space?
  \item \textbf{Comparison:} Compare file sizes — a text document, a photo, a song, and a video. Why are they so different?
  \item \textbf{Analogy building:} Ask children to invent their own analogy for RAM vs storage (not a desk — something from their own life).
\end{itemize}

\subsection*{Curriculum Links}
Mathematics (units: KB, MB, GB, TB), Computing (data representation), Science (capacity).

\sectionrule

% ============================================================
\section*{Chapter 4: The Amazing Touchscreen}
% ============================================================

\subsection*{Discussion Questions}
\begin{enumerate}
  \item Why does a touchscreen not work when you use a gloved hand?
  \item Why do some apps respond differently to a swipe versus a tap?
  \item Can you think of other surfaces in everyday life that respond to touch?
\end{enumerate}

\subsection*{Extension Activities}
\begin{itemize}
  \item \textbf{Glove test:} Try using a touchscreen with a rubber glove, then a knitted glove, then a bare finger. Discuss why each behaves differently.
  \item \textbf{Gesture dictionary:} List and define all the touch gestures you know (tap, swipe, pinch, spread, long press, etc.).
  \item \textbf{Design challenge:} Invent a new touchscreen gesture and describe what it should do in an app.
\end{itemize}

\subsection*{Curriculum Links}
Science (electricity, conductivity), Computing (human-computer interaction), Art (design thinking).

\sectionrule

% ============================================================
\section*{Chapter 5: Tiny Cameras, Big Pictures}
% ============================================================

\subsection*{Discussion Questions}
\begin{enumerate}
  \item Why do phone cameras keep improving even though the phone size stays the same?
  \item What is the difference between digital zoom and optical zoom? Which produces a better image?
  \item How do you think night mode ``finds'' more light when there is not much there?
\end{enumerate}

\subsection*{Extension Activities}
\begin{itemize}
  \item \textbf{Light experiment:} Photograph the same scene in bright light, dim light, and with night mode on. Compare the results.
  \item \textbf{Pixel art:} Create pixel art on squared paper to understand how images are made of tiny coloured squares.
  \item \textbf{Pinhole camera:} Build a simple pinhole camera from a cardboard box to see how lenses focus light.
\end{itemize}

\subsection*{Curriculum Links}
Science (light, optics), Art (photography, composition), Computing (digital images, data).

\sectionrule

% ============================================================
\section*{Chapter 6: The Battery Power Station}
% ============================================================

\subsection*{Discussion Questions}
\begin{enumerate}
  \item Why do batteries wear out over time and need replacing after a few years?
  \item Which things use the most battery on a phone? How could you test this?
  \item Where does the electricity that charges your phone originally come from?
\end{enumerate}

\subsection*{Extension Activities}
\begin{itemize}
  \item \textbf{Battery diary:} Track battery percentage at set times throughout the day. Which activities drain it fastest?
  \item \textbf{Lemon battery:} Build a simple lemon battery to demonstrate that chemical reactions can produce electricity.
  \item \textbf{Energy chain:} Trace the energy chain from a power station to the phone screen.
\end{itemize}

\subsection*{Curriculum Links}
Science (energy transfer, chemical energy, electricity), Geography (energy sources), PSHE (environmental impact).

\sectionrule

% ============================================================
\section*{Chapter 7: The Hidden Sensors}
% ============================================================

\subsection*{Discussion Questions}
\begin{enumerate}
  \item Why does your screen turn off automatically when you hold the phone to your ear? Which sensor does that?
  \item How do fitness apps know how many steps you have taken without you pressing anything?
  \item If you were designing a new sensor for smartphones, what would it detect?
\end{enumerate}

\subsection*{Extension Activities}
\begin{itemize}
  \item \textbf{Sensor hunt:} Use a free sensor app to display readings from the accelerometer and gyroscope in real time. Tilt, shake, and spin the phone.
  \item \textbf{Compass challenge:} Use the compass app to navigate from one room to another using only compass bearings.
  \item \textbf{Mapping sensors:} Draw a phone outline and annotate where each sensor sits and what it does.
\end{itemize}

\subsection*{Curriculum Links}
Science (forces, magnetism, sound, light), Computing (input devices), Mathematics (measurement, direction).

\sectionrule

% ============================================================
\section*{Chapter 8: Apps: The Workers Inside Your Phone}
% ============================================================

\subsection*{Discussion Questions}
\begin{enumerate}
  \item What is the difference between a free app and a paid app? Are free apps really free?
  \item Why is it important to read what permissions an app asks for before installing it?
  \item What would a world without apps look like? What would you miss most?
\end{enumerate}

\subsection*{Extension Activities}
\begin{itemize}
  \item \textbf{App audit:} List all apps on a device and categorise them by purpose. Which category has the most?
  \item \textbf{Permission check:} Look at the permissions a familiar app requests. Discuss whether each permission makes sense for what the app does.
  \item \textbf{Design your own app:} Sketch a paper prototype for an app that solves a real problem in your life.
\end{itemize}

\subsection*{Curriculum Links}
Computing (software design, algorithms), PSHE (privacy, digital literacy), Citizenship (consumer awareness).

\sectionrule

% ============================================================
\section*{Chapter 9: GPS: Finding Your Way}
% ============================================================

\subsection*{Discussion Questions}
\begin{enumerate}
  \item GPS works in remote forests and deserts. Why do you think it works even with no phone signal?
  \item How did people navigate before GPS? What tools did they use?
  \item Can you think of any problems with always knowing where you are (or others knowing where you are)?
\end{enumerate}

\subsection*{Extension Activities}
\begin{itemize}
  \item \textbf{Satellite map:} Look up how many GPS satellites are currently in orbit. Draw their orbits around a globe.
  \item \textbf{Triangulation activity:} Use three circles drawn on a map (each representing the range from a different point) to find where they overlap — that is triangulation.
  \item \textbf{Paper map navigation:} Navigate a short local route using only a printed paper map — no phone allowed.
\end{itemize}

\subsection*{Curriculum Links}
Geography (navigation, mapping, coordinates), Science (satellites, orbits), Mathematics (geometry, measurement).

\sectionrule

% ============================================================
\section*{Chapter 10: Stay Safe and Smart}
% ============================================================

\subsection*{Discussion Questions}
\begin{enumerate}
  \item Why is a fingerprint or face unlock more secure than a simple number PIN?
  \item What personal information should you \emph{never} share in an app or online?
  \item What does ``being kind online'' mean to you? Is it different from being kind in person?
\end{enumerate}

\subsection*{Extension Activities}
\begin{itemize}
  \item \textbf{Password strength test:} Use an offline password checker to rate the strength of sample passwords (not real ones).
  \item \textbf{Permission roleplay:} Act out scenarios where an app asks for a suspicious permission — children decide whether to allow or deny.
  \item \textbf{Digital kindness pledge:} Create a class or family digital kindness charter and display it near your devices.
\end{itemize}

\subsection*{Curriculum Links}
PSHE (online safety, digital wellbeing), Computing (cybersecurity basics), Citizenship (rights and responsibilities online).

\sectionrule

% ============================================================
\section*{Connections to Book 1}
% ============================================================

\begin{tcolorbox}[enhanced,colback=TechPurple!5,colframe=TechPurple!50,arc=6pt,boxrule=1pt,
  left=10pt,right=10pt,top=10pt,bottom=10pt]
\begin{itemize}
  \item \textbf{Chapter 1} (What Is a Smartphone?) links to \textbf{Book 1, Chapter 1} (What Is the Internet?) — both are entry points into digital technology.
  \item \textbf{Chapter 8} (Apps) links to \textbf{Book 1, Chapter 2} (Websites and Browsers) — both explain how software is the layer users interact with.
  \item \textbf{Chapter 9} (GPS) extends \textbf{Book 1, Chapter 6} (How Messages and Videos Travel) — data packets vs.\ satellite signals.
  \item \textbf{Chapter 10} (Stay Safe and Smart) directly continues \textbf{Book 1, Chapter 9} (Stay Safe and Curious) with device-specific safety.
  \item \textbf{Consider adding}: A connectivity chapter (Wi-Fi, Bluetooth, 4G/5G) to close the loop with Book 1's Wi-Fi and routers content.
\end{itemize}
\end{tcolorbox}

\vspace{12pt}

\begin{tcolorbox}[enhanced,colback=LightGray,colframe=TechPurple!40,arc=6pt,boxrule=0.8pt,
  left=10pt,right=10pt,top=8pt,bottom=8pt]
\textbf{Series:} The Secret Life of Technology\\
\textbf{Book 1:} The Secret Life of the Internet \checkmark\\
\textbf{Book 2:} Inside a Smartphone: The Tiny World in Your Hand \checkmark
\end{tcolorbox}
